\documentclass[10pt]{article}
\usepackage{icmlko}

\icmlmeta{IP 파운데이션 월드모델}{972}{크기는 못 가르고 부호는 가른다}{같은 유보 예측에 진짜 귀무 하나를 걸었더니 크기 통계량은 $p = 0.102$ 로 떨어지고 부호 일관성은 $p = 0.009$ 로 넘는다 --- 그리고 앞 노트의 자 계수는 커밋본에서 재현되지 않았다}

\begin{document}

\twocolumn[
\icmlheader

\begin{abstract}
앞 노트(971)는 아홉 사이클 만에 자를 \emph{연관}에서 \emph{유보 예측}으로 옮기고
$\Delta\rho_{\mathrm{pred}} = +0.025527$ 을 냈다. 그 값은 재현된다 --- 이 노트가 같은 생산
함수를 불러 \textbf{같은 소수 여섯 자리}를 다시 냈다. 문제는 값이 아니라 \textbf{판정 규칙}이었다.
971 의 채택 조건 다섯 중 셋(①②③)은 \textbf{같은 스칼라의 포개진 문턱}이고,
다섯째(「7 중 $\ge 5$ 도메인 양수」)는 \textbf{귀무에서 22.4\% 로 그냥 통과}한다 ---
표본분포를 가진 검정은 사실상 하나뿐이었다.
이 노트는 채택 조건을 \textbf{진짜 귀무 하나}로 다시 사전등록했다: 도메인 안에서 유보 결과를
섞는 순열 귀무(4000 뽑기 · 중심 $-1\times10^{-5}$). 결과는
\textbf{$p = 0.102474$ --- 못 넘었다.} 그러나 \textbf{같은 귀무에 부호 일관성}(7 도메인 전부 양수)을
걸면 \textbf{$p = 0.008998$ --- 넘는다.} 🔴 \textbf{한 자료에서 크기는 잡음과 안 갈리고 방향은 갈린다.}
폭도 다시 냈다: 971 의 붓스트랩은 유보만 되뽑았고(SD $0.020436$), \textbf{학습까지 되뽑으면}
$0.012912$, \textbf{이중이면 $0.027010$} 이다. 묶음 $\Delta$ 의 \textbf{42.8\%} 가 학습 41 행짜리
한 도메인에서 오고, 그 하나를 빼면 $0.015846$ 으로 \textbf{자기 잡음 바닥 $0.017991$ 아래}다.
곁으로, 앞 노트가 낸 자 계수는 \textbf{커밋된 트리에서 재현되지 않는다} ---
자 팔이 돈 뒤 대상 파일이 고쳐졌고 산출물을 다시 안 냈다. 커밋본에서 다시 재니
\textbf{증명된 낙하가 4 가 아니라 2} 이고 \textbf{밑값보다 낫지 않다.}
\end{abstract}
]

\section{왜 조건 다섯을 버렸나}

앞 노트는 측정 \emph{전}에 채택 조건 다섯을 못 박았다. 그 규율 자체는 옳다.
그런데 다섯이 \textbf{독립된 다섯 가지 물음이 아니었다.}

\begin{itemize}\itemsep2pt
\item ① $\Delta > 0$ · ② $\Delta \ge 0.01055$ · ③ $\Delta \ge$ 팔 R 95 분위 ---
      셋은 \textbf{같은 스칼라 $\Delta$ 에 대한 포개진 문턱}이다. ③이 참이면 ①②는 거의 자동이다.
\item ⑤ 「7 중 $\ge 5$ 도메인 양수」 --- 이 노트가 순열 귀무에서 재니
      \textbf{22.4\%} 의 뽑기가 그냥 통과한다. \textbf{자가 아니라 문턱이 헐거웠다.}
\item ④ 짝 붓스트랩 아래끝 --- \textbf{표본분포를 가진 것은 이것뿐이었고 그것이 떨어졌다.}
\end{itemize}

\claimx{채택 조건을 여럿 세어 「n/5」로 적으면 그 수는 검정력이 아니라 조건의 상관을 센 것이다}{조건들이 서로 독립이면 --- 즉 하나를 못 넘을 때 나머지가 같이 안 흔들리면 --- 이 주장은 떨어진다. 여기서는 ①②③이 한 스칼라의 포개진 문턱이고 ⑤가 귀무에서 22.4\% 로 통과한다.}

그래서 972 는 조건을 \textbf{하나}로 줄였다. 도메인 \emph{안에서} 유보 결과 $y_{\mathrm{ho}}$ 를
섞는다. 예측 $p_B$ · $p_C$ 는 \textbf{학습 행에서만} 적합돼 고정이므로, 섞기는
\textbf{결과와 예측의 짝만} 깬다. 앞 노트가 곁으로 쓴 「무작위 열」 귀무는 중심이 음수라
반보수적이었는데, 이 귀무는 \textbf{중심이 $-1\times10^{-5}$}(4000 뽑기)로 0 에 붙는다.

\section{크기와 부호가 갈린다}

\begin{center}\tsize
\begin{tabular}{lrr}
\toprule
통계량 & 값 & 순열 $p$ \\
\midrule
묶음 $\Delta\rho_{\mathrm{pred}}$ & $+0.025527$ & $\mathbf{0.102474}$ \\
7 도메인 전부 양수 & $7/7$ & $\mathbf{0.008998}$ \\
\midrule
\multicolumn{2}{l}{971 의 조건 ⑤($\ge 5$)가 귀무에서 통과} & $22.4\%$ \\
\bottomrule
\end{tabular}
\end{center}

같은 자료 · 같은 귀무 · 같은 4000 뽑기다. \textbf{답이 갈린다.}
크기 통계량은 유보 행으로 가중하므로 \textbf{세계애니(유보 230 행 · $\Delta = +0.002545$)}
같은 큰 도메인이 묶음을 누른다. 부호 통계량은 크기를 안 보고 \textbf{방향만} 본다.

\claimx{이 자료에서 들뜸이 유보 예측을 돕는다는 증거는 크기가 아니라 부호 일관성에 있다}{같은 순열 귀무에서 두 $p$ 가 같은 쪽으로 떨어지면 --- 즉 부호 일관성도 0.05 를 못 넘으면 --- 이 주장은 떨어진다.}

\section{폭 --- 학습을 되뽑으면}

971 의 짝 붓스트랩은 \textbf{유보 행만} 되뽑았다. 그 폭은 「적합을 고정했을 때」의 것이다.
972 는 셋을 다 냈다(2000 뽑기 · 씨앗 972).

\begin{center}\tsize
\begin{tabular}{lrr}
\toprule
되뽑는 것 & 짝 SD & 95\% 구간 \\
\midrule
유보만(971 판) & $0.020436$ & $[-0.0155,\ 0.0639]$ \\
학습만(다시 적합) & $0.012912$ & $[-0.0019,\ 0.0498]$ \\
\textbf{이중} & $\mathbf{0.027010}$ & $[-0.0237,\ 0.0843]$ \\
\bottomrule
\end{tabular}
\end{center}

\textbf{셋 다 0 을 품는다.} 학습만 되뽑은 폭이 유보만보다 \emph{좁다}는 것도 수다 ---
이 표본에서 변동의 큰 몫은 \textbf{적합이 아니라 채점 행}에서 온다.

\section{잎 하나가 절반 가까이 든다}

\begin{center}\tsize
\begin{tabular}{lrrr}
\toprule
도메인 & 유보 & 학습 & 묶음 $\Delta$ 기여 \\
\midrule
\textbf{시장팝업} & $42$ & $41$ & $\mathbf{42.8\%}$ \\
애니 & $114$ & $248$ & $32.9\%$ \\
웹툰 & $38$ & $158$ & $8.5\%$ \\
도서 & $27$ & $20$ & $5.7\%$ \\
게임 & $52$ & $93$ & $5.4\%$ \\
세계애니 & $230$ & $845$ & $4.3\%$ \\
모바일 & $29$ & $120$ & $0.4\%$ \\
\bottomrule
\end{tabular}
\end{center}

학습 \textbf{41 행}짜리 도메인 하나가 묶음의 \textbf{42.8\%} 를 든다.
그 하나를 빼면 묶음 $\Delta$ 가 $\mathbf{0.015846}$ 으로 내려가고,
이 자로 직접 잰 잡음 바닥(무작위 열 \textbf{2000} 뽑기의 95 분위) $\mathbf{0.017991}$ \textbf{아래}다.
🔴 앞 노트가 쓴 바닥 $0.016593$ 은 \textbf{50 뽑기}의 값이고, 972 가 같은 씨앗의 첫 50 만
다시 뽑아 \textbf{비트로 재현}했다 --- \textbf{바닥이 뽑기 수에 따라 움직인다}는 뜻이다.

\section{앞 노트의 두 수를 정정한다}

\subsection{「연관이 예측보다 5.7 배」 --- 정정}

앞 노트는 유보 행 가중 $|\rho_{\text{연관}}| = 0.145443$ 을 $\Delta = 0.025527$ 로 나눠
\textbf{5.7 배}라 적었다(\emph{정정 대상}). 🔴 그 절댓값이 수를 만들었다.
\textbf{부호를 살리면} 유보 행 가중 연관은 $\mathbf{+0.027166}$ 이고 배수는 $\mathbf{1.06}$ 배다.
유보 연관이 \textbf{음수}인 도메인이 셋이고 그것이 유보 행의 \textbf{56.2\%} 를 차지한다.

\subsection{「부호가 뒤집히는 셋」 --- 정정}

앞 노트는 \emph{도서 · 세계애니 · 시장팝업}에서 「학습에서 배운 부호가 유보에서 뒤집힌다」고
적었다(\emph{정정 대상}). 그 판정은 $\rho_A$(들뜸만으로 낸 \emph{유보 예측})와 유보 연관을 견준 것이라
\textbf{원리상 뒤집힘을 못 본다} --- 이 자료에서 $|\rho_A| = |\rho_{\text{유보 연관}}|$ 이 7/7 이고
$\rho_A > 0$ 도 7/7 이다. 🔴 \textbf{훈련 행 안의 연관을 내는 커밋된 러너가 0 개였다.}
972 가 그것을 내니 부호가 갈리는 도메인은 \textbf{모바일 · 시장팝업 · 웹툰}(유보 행 20.5\%)이고
앞 목록과 \textbf{하나만 겹친다}.

\claimx{같은 유보 행에서 잰 연관과 예측 이득은 절댓값을 떼면 거의 같은 크기다}{부호를 살린 유보 가중 연관을 예측 $\Delta$ 로 나눈 값이 2 를 넘으면 이 주장은 떨어진다. 여기서는 1.06 이다.}

\section{자 팔이 커밋본에서 재현되지 않았다}

앞 노트의 자 산출물은 대상 파일의 sha 를 \texttt{7a0f3f65\ldots} 로 적었는데
커밋된 blob 은 \texttt{c870fe5d\ldots} 다 --- \textbf{자 팔이 돈 뒤 그 파일이 고쳐졌고
산출물을 다시 안 냈다.} 972 가 커밋본에서 같은 자를 다시 걸었다.

\begin{center}\tsize
\begin{tabular}{lrrr}
\toprule
자 판 & 분모 & 증명된 낙하 & \% \\
\midrule
971 이 적은 것 & 13 & 4 & 30.8 \\
\textbf{커밋본(접미 일치)} & 13 & \textbf{2} & \textbf{15.4} \\
커밋본(정확 일치) & 12 & 2 & 16.7 \\
밑값 \texttt{recommit970}(정확) & 6 & 1 & 16.7 \\
\bottomrule
\end{tabular}
\end{center}

🔴 \textbf{「밑값을 넘었다」가 「밑값과 같거나 못하다」로 뒤집힌다.}
그리고 자기 자신에게 걸었더니 \textbf{자의 §4 검사 하나가 항진명제로 잡힌다} ---
여섯 판(생성기 판 $\times$ 키 일치 규칙 $\times$ 풀 범위) \textbf{전부에서} 그렇다.

\section{자가 파일의 나머지 내용에 매여 있었다}

자 B 는 무작위 입력을 만들어 검사가 상수인지 본다. 그 \emph{생성기 풀}이
\textbf{파일 전체 AST} 의 문자열과 딕트 열쇠에서 만들어졌다. 972 는 풀을
\textbf{그 자리를 감싸는 함수}로 좁혔고(기본값을 바꿨다), 둘을 나란히 냈다.
같은 커밋본에 대해 \textbf{범위만 바꿔도 낙하가 2 에서 3 으로 갈린다.}

그리고 \textbf{주장 자체를 실측했다.} 대상 파일 셋에
\textbf{검사 자리를 하나도 안 더하고 기존 자리의 줄 번호도 안 바꾸는} 꼬리를
갈래마다 12 판씩 붙여, 판정이 뒤집히는 \emph{판의 비율}을 셌다.

\begin{center}\tsize
\begin{tabular}{lrr}
\toprule
꼬리 갈래 & \texttt{file} 범위 & \texttt{func} 범위 \\
\midrule
48 줄(무관) & $\mathbf{12/36}$ & $\mathbf{0/36}$ \\
풀 잘림을 넘김 & $\mathbf{24/36}$ & $\mathbf{0/36}$ \\
주석만(음성 대조) & $0/36$ & $0/36$ \\
\bottomrule
\end{tabular}
\end{center}

🔴 \textbf{판정과 무관한 48 줄이 옛 범위에서는 셋에 하나 꼴로 판정을 뒤집고, 새 범위에서는
한 번도 안 뒤집는다.} 주석만 붙인 음성 대조는 양쪽 다 0 이다 --- 자가 AST 만 본다는 뜻이다.
🔴 \textbf{그러므로 이 저장소가 965 이후 낸 모든 「항진명제 $n$ · 낙하 $n$」은
그 파일의 나머지 내용에 매인 수다.}

\claimx{자기 검사 계수를 인용할 때는 생성기 판·키 일치 규칙·풀 범위를 같이 적어야 한다}{셋 중 어느 것을 바꿔도 같은 파일의 같은 자리에서 같은 계수가 나오면 이 주장은 떨어진다. 여기서는 풀 범위만 바꿔도 낙하가 2 와 3 으로 갈린다.}

\section{배선이 이번엔 죽지 않았다}

앞 노트의 배선 아홉 자리는 \textbf{생산 함수 일곱을 동시에 부숴도 전부 초록}이었다 ---
자리들이 생산 함수를 부르지 않고 \emph{안에서 다시 짰기} 때문이다.
972 의 배선은 전부 \texttt{predict971} 의 생산 함수를 \textbf{부른다}. 그것을 증명하려고
파괴 대조를 넣었다: 생산 함수를 \textbf{하나씩} 상수로 망가뜨리고 붉어지는 자리를 센다 ---
\textbf{9 / 9}, 즉 \textbf{하나도 빠짐없이} 잡힌다.
누출 검사도 \textbf{양방향}으로 바꿨다(유보를 갈면 안 변해야 하고 \textbf{학습을 갈면 변해야 한다}) ---
앞 판은 적합기를 상수 0 으로 망가뜨려도 초록이었다.
자기 참조 검사 하나(리터럴을 인자로 넘겨 같은 상수와 비교하던 자리)는 \textbf{지웠다}.

\section{남은 것}

\begin{itemize}\itemsep2pt
\item \textbf{크기를 가르려면 표본이 더 필요하다.} 이 표본의 묶음 최소검출효과(짝 안 지음 상한)는
      $\mathbf{0.087292}$ 다 --- 앞 노트가 「약 0.090」이라 적은 수의 실값이다(\emph{정정}).
\item \textbf{부호 일관성의 분모가 7 이다.} 도메인을 쪼개거나 창을 갈라 분모를 늘려야 한다.
      학습 문턱을 판과 같은 15 로 낮추면 도메인이 \textbf{8} 이 되고 묶음이 $0.020581$ 이다(곁 팔).
\item \textbf{「긴 띠가 정보를 갖는다」와 「길수록 좋다」가 아직 안 갈렸다.} 들뜸은 차 하나라
      길이-반응 곡선이 없다.
\end{itemize}

\end{document}
